% =====================================================================
%  INTRODUCCIÓN  —  párrafos continuos, sin subsecciones (lo pide la plantilla)
%  Fuentes: anteproyecto cap. 1 (introducción) y cap. 3 (justificación).
%  Registro: frases cortas, una idea por frase. Mediana objetivo ~15 palabras.
%
%  CITAS: se usa \textcite, que NOMBRA al autor dentro de la frase y deja el
%  numero como localizador. La forma «segun Fulano, tal cosa» es la que pidio el
%  director: una cita que solo es un numero entre corchetes no dice de quien es
%  la idea, y obliga al lector a irse al final del documento para saberlo.
% =====================================================================
\chapter*{Introducción}
\addcontentsline{toc}{chapter}{INTRODUCCIÓN}

Orientarse dentro de un edificio grande no es igual de fácil para todo el mundo.
Quien llega por primera vez construye un mapa mental a partir de información
parcial, y lo va corrigiendo mientras camina. Ese esfuerzo crece cuando el
espacio tiene pasillos repetidos, pocos puntos de referencia o señalización
escasa. Y crece más cuando hay que cambiar de piso, porque la escalera corta la
continuidad de lo que la persona lleva recorrido. Jiménez Muñoz lo señala al construir una
arquitectura cognitiva que mete al usuario dentro de un sistema de mapas de
varios niveles: la escalera es el punto donde la representación mental del
recorrido deja de servir \autocite{jimenez2024arquitectura}.

La robótica móvil colaborativa lleva años resolviendo tareas parecidas. El
reparto de trabajo entre varios robots tiene base teórica desde que
\textcite{gerkey2004taxonomy} formalizaron la asignación de tareas en sistemas
multi-robot. Y \textcite{varghese2023multirobot} muestran que varios agentes
pueden mapear y localizarse de forma cooperativa. Casi todos esos sistemas
operan, sin embargo, en una sola planta.

El cambio de nivel se ha abordado por dos caminos, y los dos pasan por el
hardware. \textcite{kim2024indoor} combinan un mapa por piso con un grafo
topológico y hacen que el robot use el ascensor. \textcite{song2024elevator} van
más lejos: automatizan ese uso instalando dispositivos conectados en el
edificio. La otra opción es dotar al robot de mecanismos para subir escaleras.
Ambas encarecen la plataforma y exigen intervenir la infraestructura.

Este trabajo propone otra. En vez de un robot que cruza de piso, se usan \textbf{dos agentes}, robots móviles dedicados cada uno a un nivel. El usuario sube por su cuenta y el segundo agente lo recibe arriba. El cambio de piso deja de ser un problema
de locomoción y pasa a ser uno de coordinación. Dos agentes se reparten una tarea
y se la pasan en un punto acordado. Ese traspaso es el \emph{relevo}, y es el aporte que el proyecto declara.

La pregunta que organiza el trabajo es cómo asistir la orientación de usuarios
dentro de entornos interiores con múltiples pisos utilizando sistemas robóticos,
y el relevo es la respuesta que se propone y se pone a prueba. No basta con que los agentes se muevan. Hay que demostrar que la asistencia no se interrumpe cuando
el usuario cambia de nivel, y medirlo con criterios fijados antes de correr los
experimentos.

Los agentes se construyeron sobre la plataforma AWS DeepRacer, con software libre. La
coordinación corre en un nodo central que habla con cada agente a través de
ROS~2. \textcite{macenski2022ros2} lo describen como un rediseño sobre DDS, no
como una actualización de su antecesor. El usuario pide el
servicio desde una página web que funciona en cualquier teléfono, sin instalar
nada y sin salida a internet.
La evaluación se hizo primero en simulación, con una campaña de treinta misiones
sorteadas, y después con los agentes físicos, con una campaña de diez misiones sorteadas entre los pisos 3 y 4 del edificio Calatayud de la universidad.

